\section{Заключение}

В ходе выполнения лабораторной работы были изучены основные принципы
организации разработки программного обеспечения с использованием
методологии Scrum и получены практические навыки управления проектом
в Jira.

В рамках работы был создан Scrum-проект, сформирован состав команды и
распределены роли участников. Была создана и приоритизирована очередь
продукта, выполнена оценка пользовательских историй, сформирован выпуск
и определён состав первого спринта.

В процессе планирования спринта пользовательские истории были
декомпозированы на отдельные задачи, для которых были заданы оценки
трудоёмкости и назначены исполнители. В ходе моделирования выполнения
спринта изменялись состояния задач, фиксировался прогресс работы и
отслеживался оставшийся объём работ.

Для контроля состояния проекта были использованы запросы, отчёты и
сводная панель Jira. Также были рассмотрены процедуры проведения
Sprint Review и ретроспективы, завершения спринта и переноса
незавершённых задач в следующую итерацию.

Таким образом, в результате лабораторной работы были освоены основные
этапы управления программным проектом по Scrum: от формирования
Product Backlog и планирования выпуска до организации спринта,
контроля выполнения работ и подготовки следующей итерации.
